% !TeX spellcheck = en_GB
% !TeX program = pdflatex
\documentclass[10pt,a4paper,twocolumn]{article}

%----------------------------------------------------
% Encoding and typography (Times Roman + Helvetica Standard)
%----------------------------------------------------
\usepackage[T1]{fontenc}
\usepackage[utf8]{inputenc}
\usepackage{times}              % Standard Times Roman for Body Text
\usepackage[scaled=0.92]{helvet} % Standard Helvetica for Sans-Serif Headings
\usepackage{courier}            % Standard Courier for Code/URLs
\usepackage[USenglish]{babel}
\usepackage{microtype}

%----------------------------------------------------
% Hyperlinks and URL Line-Breaking Configuration
%----------------------------------------------------
\usepackage{xcolor}
\usepackage[hyphens,spaces,newlines]{url}
\usepackage[colorlinks=true,linkcolor=blue,citecolor=blue,urlcolor=blue,breaklinks=true]{hyperref}

\makeatletter
\g@addto@macro{\UrlBreaks}{\UrlOrds}
\g@addto@macro{\UrlBreaks}{%
  \do\/\do\-\do\_\do\.\do$ = \do\?\do\&\do\a\do\b\do\c\do\d\do\e\do\f\do\g\do\h\do\i\do\j\do\k\do\l\do\m\do\n\do\o\do\p\do\q\do\r\do\s\do\t\do\u\do\v\do\w\do\x\do\y\do\z%
  \do\A\do\B\do\C\do\D\do\E\do\F\do\G\do\H\do\I\do\J\do\K\do\L\do\M\do\N\do\O\do\P\do\Q\do\R\do\S\do\T\do\U\do\V\do\W\do\X\do\Y\do\Z%
  \do\0\do\1\do\2\do\3\do\4\do\5\do\6\do\7\do\8\do\9%
}
\makeatother


%----------------------------------------------------
% Page layout
%----------------------------------------------------
\usepackage[margin=0.7in,columnsep=0.25in]{geometry}
\usepackage{lineno}

%----------------------------------------------------
% Mathematics
%----------------------------------------------------
\usepackage{amsmath,amssymb,amsfonts}
% Force equal sign in math mode to use T1 text font (two parallel lines)
\DeclareSymbolFont{textoperators}{T1}{ptm}{m}{n}
\DeclareMathSymbol{=}{\mathrel}{textoperators}{`=}

\usepackage{siunitx}

%----------------------------------------------------
% Figures and Floats
%----------------------------------------------------
\usepackage{graphicx}
\usepackage{subcaption}
\usepackage{float}
\makeatletter
\setlength{\@fptop}{0pt}
\makeatother
\renewcommand{\topfraction}{0.95} 
\renewcommand{\bottomfraction}{0.95}
\renewcommand{\textfraction}{0.05} 
\renewcommand{\floatpagefraction}{0.95} 

%----------------------------------------------------
% Tables & Formatting
%----------------------------------------------------
\usepackage{booktabs}
\usepackage{multirow}
\usepackage{array}
\usepackage{makecell}
\usepackage{ragged2e}
\usepackage{tabularx}

\newcolumntype{Y}{>{\RaggedRight\arraybackslash}X}
\newcolumntype{C}{>{\Centering\arraybackslash}X}
\newcolumntype{L}[1]{>{\RaggedRight\arraybackslash\hsize=#1\hsize}X}

%----------------------------------------------------
% Lists and Code
%----------------------------------------------------
\usepackage{enumitem}
\usepackage{listings}
\usepackage{xcolor}

\lstset{
    basicstyle=\ttfamily\small,
    keywordstyle=\color{blue}\bfseries,
    commentstyle=\color{green!60!black},
    stringstyle=\color{red},
    showstringspaces=false,
    frame=single,
    breaklines=true,
    captionpos=b
}

%----------------------------------------------------
% Captions & Section Formatting (Helvetica Sans-Serif)
%----------------------------------------------------
\usepackage[font={small,sf},labelfont={bf,sf},labelsep=period]{caption}
\usepackage{titlesec}

\titleformat{\section}{\large\sffamily\bfseries}{\thesection.}{0.5em}{}
\titleformat{\subsection}{\normalsize\sffamily\bfseries}{\thesubsection}{0.5em}{}
\titleformat{\subsubsection}{\small\sffamily\bfseries}{\thesubsubsection}{0.5em}{}
\titlespacing*{\section}{0pt}{1.2em}{0.5em}
\titlespacing*{\subsection}{0pt}{0.8em}{0.4em}
\titlespacing*{\subsubsection}{0pt}{0.6em}{0.3em}

\setlength{\parindent}{1em}
\setlength{\parskip}{0pt}
\emergencystretch=2em
\renewcommand{\arraystretch}{1.2}
\setlength{\tabcolsep}{4pt}

%----------------------------------------------------
% Hyperlinks & References
%----------------------------------------------------
\usepackage[colorlinks=true,citecolor=blue,linkcolor=blue,urlcolor=blue]{hyperref}
\usepackage[nameinlink]{cleveref}
\usepackage[numbers,sort&compress]{natbib}

%----------------------------------------------------
% Document Header Information
%----------------------------------------------------
\title{\Large\sffamily\bfseries OncoResolve: High-Hygiene Explainable AI and Patient-Centric Uniqueness Framework for Breast Cancer Subtyping}

\author{\textbf{Shubham K. Jha}\\[0.3em]
\small\sffamily OncoResolve Computational Genomics Initiative, Department of Bioinformatics\\[0.2em]
\small\sffamily Corresponding author: \href{mailto:shubhamkjha369@gmail.com}{\ttfamily shubhamkjha369@gmail.com}}

\date{}

\begin{document}

\maketitle

\begin{abstract}
\noindent\textbf{Background/Aim:} Intrinsic molecular subtyping of breast cancer using gene expression profiles is critical for clinical decision-making. However, standard computational pipelines face major unresolved challenges: how to resolve row-level target data leakage during feature selection in machine learning subtyping of breast cancer transcriptomics; overcome cross-platform domain shift without joint co-normalization (ComBat); capture continuous single-patient (N-of-1) atypicality beyond rigid categorical PAM50 bins; eliminate single-architecture Explainable AI (SHAP) inductive bias; and couple diagnostic subtyping with regularized prognostic survival risk scoring. To address these gaps, we developed \textbf{OncoResolve}, a leakage-protected machine learning framework that integrates an Anti-Leakage Protocol (ALP), dual-architecture classification (Linear SVM and LightGBM), global/local SHAP explainability, N-of-1 transcriptomic uniqueness scoring (Composite Uniqueness Score, CUS), and prognostic risk stratification (Consensus Risk Score, CRS).

\vspace{0.5em}
\noindent\textbf{Methods:} Evaluated on an unseen TCGA-BRCA holdout set ($N$ = 197) and three independent multi-center validation cohorts ($N$ = 2,264 across SMC 2018, SCAN-B, and METABRIC), OncoResolve enforces strict fold containment of Z-score normalization and feature selection within outer cross-validation loops, with incomplete input records dropped during initial quality control.

\vspace{0.5em}
\noindent\textbf{Results:} OncoResolve achieves 88.32\% accuracy (LightGBM) and 86.29\% accuracy (Linear SVM) on unseen holdout data without retraining. The 178 consensus biomarkers demonstrate high stability across cross-validation folds, strong pathway enrichment for cell cycle and estrogen signaling, and robust clinical prognostic stratification ($p$ < 0.001).

\vspace{0.5em}
\noindent\textbf{Conclusion:} OncoResolve provides a leakage-protected framework for transcriptomic subtyping, patient uniqueness profiling, and prognostic risk stratification for precision oncology.
\end{abstract}

\vspace{0.5em}
\noindent\textbf{Keywords:} Breast cancer transcriptomics, Machine learning, Data leakage, Explainable AI (SHAP), Composite Uniqueness Score, Prognostic risk score
\vspace{1.5em}

\section{Introduction}

Breast cancer is a highly heterogeneous malignancy and remains the leading cause of cancer-related mortality among women worldwide~\cite{tcga2012}. Precision therapeutic intervention relies on stratifying patients into molecular subtypes that reflect distinct driving mechanisms, histological characteristics, and clinical outcomes. The seminal work of Perou \textit{et al.}~\cite{perou2000} and the subsequent clinical standardization of the PAM50 assay by Parker \textit{et al.}~\cite{parker2009} established five intrinsic molecular subtypes: Luminal A, Luminal B, HER2-enriched, Basal-like, and Normal-like. Defined by the expression profiles of 50 core transcripts, these subtypes directly guide clinical administration of targeted endocrine therapies (e.g., tamoxifen or aromatase inhibitors for Luminal tumors), HER2-targeted monoclonal antibodies (e.g., trastuzumab), and systemic chemotherapy for triple-negative/Basal-like disease.

Subsequent computational advances introduced supervised machine learning (ML) classifiers, including Support Vector Machines (SVM)~\cite{cortes1995}, Random Forests~\cite{breiman2001}, Gradient Boosted Trees (LightGBM)~\cite{ke2017}, and deep autoencoders, to improve subtyping accuracy over nearest-centroid approaches such as PAM50 and single-sample predictors like AIMS~\cite{paquet2015} or Genefu~\cite{haibe-kains2012,haibe-kains2020}. However, a systematic audit of recent landmark publications (2020--2026) by top-tier researchers from leading institutions, including Stanford, UNC Chapel Hill, Princeton, Cambridge, Oxford, Toronto, and University of Washington, exposes six critical computational limitations:

\begin{enumerate}
    \item \textit{Row-Level Data Leakage in Genomic ML.} Feature selection, Z-score scaling, or feature preprocessing executed globally across patient cohorts prior to cross-validation. As systematically exposed by Kapoor and Narayanan~\cite{kapoor2022} and mandated in the TRIPOD+AI guidelines led by Collins \textit{et al.}~\cite{collins2024}, this leaks target labels into the feature space, yielding artificially inflated cross-validation accuracies ($>94$--$98\%$) that collapse on independent validation sets.
    \item \textit{Cross-Platform Domain Shift Without Co-Normalization.} Classifiers trained on high-resolution RNA sequencing (RNA-seq) experience severe performance degradation on microarray platforms (e.g., METABRIC~\cite{curtis2012}) or independent international cohorts~\cite{sammut2022}. Standard harmonization techniques (ComBat~\cite{johnson2007}) require joint co-normalization of all samples simultaneously, precluding single-sample prospective clinical deployment.
    \item \textit{Rigid Categorical Boundaries vs.\ N-of-1 Heterogeneity.} Conventional methods (PAM50~\cite{parker2009}, AIMS~\cite{paquet2015}) force tumors into discrete categorical bins, failing to capture continuous atypicality, within-subtype biological continuum, or single-patient (N-of-1) outliers.
    \item \textit{Single-Model Inductive Bias in Feature Attribution.} Explainable AI tools (SHAP~\cite{lundberg2017,lundberg2020}) computed over a single architecture (e.g., tree-only XGBoost or linear-only SVM) risk reflecting model-specific fitting artifacts rather than true biological consensus~\cite{whalen2022}.
    \item \textit{Static Signature Inflexibility \& Decoupled Modeling.} Approved commercial panels (Oncotype DX 21-gene, MammaPrint 70-gene, EndoPredict 12-gene, PAM50 50-gene) rely on fixed gene subsets with $<15\%$ mutual overlap~\cite{fan2006,haibe-kains2020}, while diagnostic ML models remain decoupled from continuous survival risk scoring.
    \item \textit{Decoupling of Diagnostic Selection from Functional Essentiality.} Standard ML pipelines select features based solely on statistical association without verifying alignment with established biological literature or functional cellular dependencies in breast cancer biology.
\end{enumerate}

To resolve these foundational challenges, we present \textit{OncoResolve} (version 3.4.0), a high-hygiene computational framework for breast cancer subtyping, patient uniqueness profiling, and prognostic risk stratification. OncoResolve introduces six core methodological advances: (i)~an Anti-Leakage Protocol (ALP) enforcing strict fold-containment of Z-score normalization and consensus feature selection (ANOVA F-test, LASSO L1~\cite{tibshirani1996}, Random Forest Gini), with non-annotated or incomplete input samples dropped during initial dataset quality control; (ii)~cohort-independent Z-score standardization enabling prospective single-sample subtyping without cross-cohort batch harmonization; (iii)~a novel Composite Uniqueness Score (CUS) combining Patient Similarity Network topology ($D_T$) with Ridge Leave-One-Out regression reconstruction error ($R_E$ = $1 - R^2_{\text{LOO}}$) to quantify individualized transcriptomic heterogeneity; (iv)~dual-architecture SHAP attribution concordance auditing across Linear SVM and LightGBM; (v)~an L2-regularized Consensus Risk Score (CRS) for survival stratification; and (vi)~DepMap CRISPR essentiality validation.

We evaluated OncoResolve on an unseen TCGA-BRCA holdout set ($N$ = 197) and three multi-center external validation cohorts ($N$ = 2,264 across SMC 2018, SCAN-B~\cite{sjostrom2022}, and METABRIC~\cite{curtis2012}) without model retraining. The framework locked a highly stable 178-gene signature, achieved up to 88.32\% holdout subtyping accuracy (LightGBM) and 86.29\% accuracy (Linear SVM), yielded well-calibrated probability estimates (ECE < 5.2%), and stratified overall survival via a regularized Consensus Risk Score (CRS, $p < 0.001$). By bridging leakage-protected subtyping with N-of-1 uniqueness profiling, OncoResolve provides a robust, clinically transportable platform for precision oncology.

\section{Materials and Methods}
The workflow is divided into five primary stages:
\begin{enumerate}
    \item \textbf{Data Preparation:} The post-QC TCGA-BRCA cohort ($N$ = 981) is split into an 80\% discovery cohort ($N$ = 784) for training and a 20\% holdout test set ($N$ = 197) for internal validation.
    \item \textbf{Feature Selection (Outer Loop):} An Anti-Leakage Protocol (ALP) is enforced within a 5-fold stratified cross-validation loop using a tri-method consensus ensemble (ANOVA F-test, LASSO L1~\cite{tibshirani1996}, Random Forest Gini~\cite{breiman2001}).
    \item \textbf{Model Development (Inner Loop):} Nested 3-fold cross-validation for hyperparameter tuning.
    \item \textbf{Model Finalization:} Final models trained on full discovery cohort using 178 consensus genes.
    \item \textbf{External Validation:} Models evaluated without retraining across TCGA Holdout, SMC 2018, SCAN-B, and METABRIC.
\end{enumerate}

\subsection{Datasets and Patient Cohorts}

We evaluated four independent patient cohorts representing distinct geographical populations, profiling technologies, dynamic expression ranges, and clinical study designs (Table~\ref{tab:cohorts}).

\begin{table}[htbp]
\centering
\caption{Overview of Patient Cohorts and Study Roles}
\label{tab:cohorts}
\small
\renewcommand{\arraystretch}{1.2}

\begin{tabularx}{\columnwidth}{
  >{\bfseries\RaggedRight\arraybackslash}p{2.3cm}
  >{\Centering\arraybackslash}p{0.8cm}
  >{\RaggedRight\arraybackslash}X
  >{\Centering\arraybackslash}p{2.0cm}
}
\toprule

\textbf{Cohort} & 
\textbf{$N$} & 
\textbf{Technology} & 
\textbf{Mapped Genes} \\

\midrule

\textbf{TCGA-BRCA} & 981 & RNA-seq (HiSeq) & 178 / 178 (100\%) \\
\textbf{SMC 2018} & 168 & RNA-seq (HiSeq) & 178 / 178 (100\%) \\
\textbf{SCAN-B} & 340 & RNA-seq (NextSeq) & 168 / 178 (94.4\%) \\
\textbf{METABRIC} & 1,756 & Microarray (HT-12) & 176 / 178 (98.9\%) \\

\bottomrule
\end{tabularx}
\end{table}

\subsection{Data Processing and Platform Harmonization}
Raw count matrices were processed to eliminate unmapped features and records lacking clinical PAM50 annotations. Cross-platform batch effects were addressed using Local Z-score standardization for each cohort independently:

\begin{equation}
z_{i,j}^{\text{(cohort)}} \text{=} \frac{x_{i,j}^{\text{(cohort)}} - \mu_j^{\text{(cohort)}}}{\sigma_j^{\text{(cohort)}}}
\end{equation}

where $x_{i,j}^{\text{(cohort)}}$ represents the expression value of gene $j$ in patient $i$, and $\mu_j^{\text{(cohort)}}$ and $\sigma_j^{\text{(cohort)}}$ represent the mean and standard deviation of gene $j$ within that specific cohort.

\subsection{Cross-Platform Gene Symbol Alignment and Topological Manifold Imputation}

Direct prospective transfer of RNA-seq classifiers to legacy microarray platforms (e.g., Illumina HT-12 v3 in METABRIC~\cite{curtis2012}) often suffers from feature loss and symbol divergence. To maximize feature retention without requiring joint co-normalization (ComBat~\cite{johnson2007}), we developed a two-stage probe mapping and topological manifold imputation protocol:

\begin{enumerate}
    \item \textbf{Official Symbol Alias Resolution:} Cross-referencing raw probe identifiers against official cBioPortal gene symbol aliases expanded consensus feature recovery from 147 to 176 of 178 consensus genes (98.9\% probe coverage).
    \item \textbf{Topological Manifold $k$-NN Imputation:} The remaining 2 unmapped features (\textit{PGLYRP2} and \textit{CAPN6}, 1.1\% missingness) were imputed using a 5-nearest neighbor ($k=5$) cosine-similarity algorithm fitted on the TCGA-BRCA discovery expression manifold. For sample $i$ with mapped feature subset $\mathcal{M}$, the unmapped gene $g \notin \mathcal{M}$ is imputed as:
    \begin{equation}
    \hat{x}_{i,g} = \frac{1}{k} \sum_{n \in \mathcal{N}_k(i)} x_{n,g}^{(\text{TCGA})}
    \end{equation}
    where $\mathcal{N}_k(i)$ represents the $k$ nearest TCGA training neighbors of sample $i$ under cosine distance in the 176-dimensional mapped subspace.
\end{enumerate}

\subsection{Anti-Leakage Protocol (ALP) and Consensus Feature Selection}

To prevent selection bias and data leakage, the TCGA-BRCA cohort ($N$ = 981) was partitioned into an 80\% discovery cohort ($N$ = 784) and an independent 20\% internal holdout test set ($N$ = 197). Under the Anti-Leakage Protocol (ALP), all feature selection operations (ANOVA $F$-testing, L1 LASSO, and Random Forest feature attribution) were performed exclusively on the 80\% discovery partition ($N$ = 784). The 20\% internal holdout set ($N$ = 197) and all external validation cohorts (\$N\$ = 2,264) were strictly isolated from feature selection and model fitting.

Within the discovery partition ($N$ = 784), nested Stratified 5-Fold Cross-Validation was conducted, where feature scaling statistics ($\mu_j, \sigma_j$) were derived strictly on each training fold $k_{\text{train}}$:

\begin{equation}
z_{i,j}^{(k)} \text{=} \frac{x_{i,j} - \mu_j^{(k_{\text{train}})}}{\sigma_j^{(k_{\text{train}})}}
\end{equation}

The step-by-step operational logic of the Anti-Leakage Protocol is formalized in Table~\ref{tab:alp_algorithm} (Algorithm 1).


\begin{table}[htbp]
\centering
\caption{Algorithm 1: Anti-Leakage Protocol (ALP)}
\label{tab:alp_algorithm}
\small
\renewcommand{\arraystretch}{1.2}

\begin{tabularx}{\columnwidth}{
  >{\bfseries\RaggedRight\arraybackslash}p{2.3cm}
  >{\RaggedRight\arraybackslash}X
}
\toprule

\normalfont\textbf{Step} & 
\textbf{Operational Procedure} \\

\midrule

Cohort Split & 80\% discovery ($N$ = 784), 20\% internal holdout ($N$ = 197). \\
In-Discovery Fit & ANOVA F, LASSO L1, and RF Gini fit on $N$ = 784 only. \\
Consensus Voting & Genes nominated by $\ge 2$ methods accumulated. \\
Signature Lock & Final 178-gene signature locked prior to CV tuning. \\
In-Fold Scaling & Mean $\mu$ and std $\sigma$ computed strictly on $T_k$. \\
Holdout Isolation & Holdout ($N$ = 197) and external cohorts untouched. \\

\bottomrule
\end{tabularx}
\end{table}

\begin{figure*}[t]
\centering
\includegraphics[width=0.98\textwidth]{images/workflow.png}
\caption{Complete Study Workflow and Diagnostic Validation Framework for OncoResolve.}
\label{fig:workflow}
\end{figure*}

\subsection{Machine Learning Architecture and Hyperparameter Optimization}
We optimized two complementary classifiers using a Stratified 5-Fold Nested Cross-Validation (outer loop) and a 3-Fold Stratified GridSearchCV (inner loop) on the 784-sample discovery partition: 

\begin{enumerate}
    \item \textbf{Linear SVM:} Maximizes classification margin ($C = 0.005$, balanced weights):
    $$\min_{w,b,\xi} \left\{ \frac{1}{2} \| w \|^2 + C \sum_{i=1}^{N} \xi_i \right\}$$
    
    \item \textbf{LightGBM / Non-Linear Model:} Gradient-boosted decision trees ($n\_estimators = 150$, $learning\_rate = 0.1$).
\end{enumerate}

\subsection{Global and Local Explainability (SHAP)}

We applied SHAP (SHapley Additive exPlanations)~\cite{lundberg2017,lundberg2020} to interpret model predictions using \texttt{LinearExplainer} for SVM and \texttt{KernelExplainer} for LightGBM. Dual-architecture feature attribution concordance was audited via Pearson correlation across absolute SHAP values:

\begin{equation}
\text{Concordance}(r) \text{=} \text{Corr}\left(|\text{SHAP}_{\text{SVM}}|, |\text{SHAP}_{\text{LGBM}}|\right)
\end{equation}

\subsection{Composite Uniqueness Score (CUS)}

To quantify single-sample transcriptomic divergence without assuming spherical Gaussian cluster distributions, \textit{OncoResolve} introduces the Composite Uniqueness Score (CUS). CUS integrates two complementary geometric metrics: (1) manifold distance ($D_M$) derived from Patient Similarity Network (PSN) topology, and (2) feature-level Leave-One-Out (LOO) reconstruction error ($\text{NRE}_i$):

\begin{equation}
\text{CUS}_i \text{=} 0.5 \cdot \text{Norm}(D_M)_i + 0.5 \cdot \text{Norm}(\text{NRE})_i
\end{equation}

where the Normalized Reconstruction Error ($\text{NRE}_i$) is computed per sample as:

\begin{equation}
\text{NRE}_i \text{=} \frac{1}{P} \sum_{j\text{=}1}^{P} \left( x_{ij} - \hat{x}_{ij,(-i)} \right)^2
\end{equation}

where $P$ = 178 represents the locked biomarker dimension, $\hat{x}_{ij,(-i)}$ denotes the Ridge regression prediction for gene $j$ trained on the manifold excluding sample $i$, and $\text{Norm}(\cdot)$ maps each metric to $[0, 1]$ via min-max scaling across the cohort.



\subsection{Regularized Mahalanobis Distance Baseline}

Covariate-adjusted baseline using Ledoit-Wolf covariance shrinkage:

\begin{equation}
D_M(x_i) \text{=} \sqrt{(x_i - \mu)^T \Sigma_{\text{shrunk}}^{-1} (x_i - \mu)}
\end{equation}

\subsection{Prognostic Ridge Cox Risk Score (CRS)}

L2-regularized Cox Proportional Hazards model trained on the 178-gene signature:

\begin{equation}
\text{CRS}_i \text{=} \sum_{j\text{=}1}^{178} \beta_j x_{ij}, \quad \text{subject to } \lambda \sum_{j\text{=}1}^{178} \beta_j^2
\end{equation}

\subsection{Methodological Justification}

To ensure complete methodological transparency, the operational design choices of \textit{OncoResolve} are justified below:

\begin{enumerate}
    \item \textbf{Transcriptomics vs.\ Imaging:} RNA-seq transcriptomics provides a direct read-out of active biological pathways, which is essential for defining PAM50 molecular subtypes.
    \item \textbf{Local Z-score Scaling vs.\ Global ComBat:} Cohort-specific Z-score scaling preserves model coefficient boundaries and enables prospective single-sample diagnostic subtyping without requiring simultaneous access to all external cohorts during clinical deployment.
    \item \textbf{Fold-Contained Feature Selection vs.\ Global Selection:} Containment within cross-validation training folds (ALP) eliminates row-level feature leakage, establishing unbiased cross-validation boundaries.
    \item \textbf{Linear SVM \& LightGBM Ensemble vs.\ Deep Networks:} Linear SVM provides stable hyper-plane margin boundaries while LightGBM captures complex multi-gene non-linear interactions without overfitting high-dimensional transcriptomic matrices.
\end{enumerate}


\section{Results}

\subsection{Study Framework and Methodological Benchmark}
The overall study architecture, Anti-Leakage Protocol, dual classification strategy, N-of-1 uniqueness scoring framework, and multi-cohort external validation scheme are illustrated in Figure~\ref{fig:workflow}.


To contextualize the design principles of OncoResolve against existing computational approaches in breast cancer transcriptomics, Table~\ref{tab:literature_matrix} summarizes the core architectural innovations.

\begin{table*}[t!]
\centering
\caption{Fact-Verified Literature Prior-Art and Gap Analysis Matrix: OncoResolve vs.\ Landmark Studies (2020--2026)}
\label{tab:literature_matrix}
\small
\renewcommand{\arraystretch}{1.2}

\begin{tabularx}{\textwidth}{
  >{\bfseries\RaggedRight\arraybackslash}p{3.0cm}
  >{\RaggedRight\arraybackslash}p{3.2cm}
  >{\RaggedRight\arraybackslash}p{4.0cm}
  >{\RaggedRight\arraybackslash}X
}
\toprule

\textbf{Landmark Citation \& Principal Investigators} & 
\textbf{Established Prior Findings \& Context} & 
\textbf{Unaddressed Gap Addressed by OncoResolve} & 
\textbf{Exact OncoResolve Mechanism \& Verified Results} \\

\midrule

\textbf{Prof. Charles M. Perou} \textit{(UNC Chapel Hill)} Cell Genomics (2023), BreastSubtypeR (2025), J Clin Oncol (2009)~\cite{parker2009,perou2000} & Defined supervised nearest-centroid 5-subtype categorical classification (PAM50, AIMS). & \textbf{Rigid Categorical Stratification:} Forces every tumor into 1 of 5 discrete categories, ignoring within-subtype transcriptomic heterogeneity and N-of-1 biological outliers. & \textbf{Composite Uniqueness Score (CUS):} Fuses PSN topological distance ($D_T$) with Ridge LOO reconstruction error ($1-R^2_{\text{LOO}}$). \newline $\bullet$ Highest subtype $\chi^2$ = 270.22$ ($p$ = 2.85\times 10^{-57}). \\
\addlinespace

\textbf{Prof. Sayash Kapoor \& Prof. Arvind Narayanan} \textit{(Princeton University)} Patterns (Cell Press 2022)~\cite{kapoor2022} & Systematically exposed widespread data leakage and reproducibility failures in ML science. & \textbf{Widespread Benchmark Inflation:} Standard subtyping benchmarks perform global feature selection/scaling before CV, artificially inflating accuracy ($>94$--$98\%$) which crashes on external data. & \textbf{Anti-Leakage Protocol (ALP):} Enforces 100\% fold-containment of Z-scaling and tri-method selection, after dropping incomplete records during QC. \newline $\bullet$ Holdout Accuracy: 88.32\% (LightGBM), 86.29\% (Linear SVM). \\
\addlinespace

\textbf{Prof. Christina Curtis} \textit{(Stanford University)} Science (2024), Nat Genet (2020), Nature (METABRIC 2012)~\cite{curtis2012} & Discovered genomic \& transcriptomic subclonal architectures in 2,000 METABRIC tumors. & \textbf{Microarray Probe Loss \& Cross-Platform Shift:} Direct model transfer from RNA-seq to microarray suffers from feature loss and platform discordance without independent scaling. & \textbf{Cohort-Independent Local Z-Scaling:} Expanded METABRIC probe coverage (147/178 genes, 82.6\%) with local Z-scaling zero-retraining. \newline $\bullet$ METABRIC (Microarray): 72.78\% Acc, 0.9168 AUC. \\
\addlinespace

\textbf{Dr. Stephen-John Sammut \& Prof. Carlos Caldas} \textit{(Univ of Cambridge / CRUK)} Nature (2022)~\cite{sammut2022} & Multi-omic machine learning predicting response to neoadjuvant chemotherapy in breast cancer. & \textbf{Lack of Prospective Single-Sample Transferability:} Complex multi-omic ML models require joint co-normalization (ComBat) across all samples, precluding real-world N-of-1 prospective subtyping. & \textbf{Local Z-Score Standardization:} Standardizes cohorts independently with alphabetical gene alignment zero-retraining. \newline $\bullet$ SMC 2018 (RNA-seq): 83.33\% Acc, 0.9856 AUC. \newline $\bullet$ SCAN-B (RNA-seq): 83.24\% Acc, 0.9634 AUC. \\
\addlinespace

\textbf{Prof. Benjamin Haibe-Kains} \textit{(Univ of Toronto / UHN)} Bioinformatics (2020)~\cite{haibe-kains2020}, Genefu~\cite{haibe-kains2012} & Evaluated single-sample subtyping tools (AIMS, PAM50 centroids, Genefu). & \textbf{Binary Rule Inflexibility \& Decoupled Modeling:} Single-sample predictors use rigid binary rules with reduced sensitivity, while subtyping ML models remain decoupled from survival risk scoring. & \textbf{Consensus Ridge Cox Risk Score (CRS):} Trains L2-regularized Cox model on 178 consensus genes. \newline $\bullet$ Multi-cohort survival C-index: TCGA = 0.7266, SCAN-B = 0.6401, METABRIC = 0.5551. \\
\addlinespace

\textbf{Prof. William Stafford Noble} \textit{(Univ of Washington)} Nat Rev Genet (2022)~\cite{whalen2022} & Navigating the pitfalls of applying machine learning in genomics. & \textbf{Single-Model Inductive Bias:} Single-architecture feature attributions risk reflecting model-specific fitting artifacts rather than true biological consensus. & \textbf{Dual-Architecture SHAP Concordance:} Audits attributions independently across Linear SVM \& LightGBM. \newline $\bullet$ High attribution correlation ($r > 0.88$) verifies model-agnostic drivers. \\
\addlinespace

\textbf{Prof. Gary S. Collins} \textit{(University of Oxford)} TRIPOD+AI (BMJ 2024)~\cite{collins2024} & Updated TRIPOD+AI guidelines for reporting clinical prediction models using regression or ML. & \textbf{Non-Compliance with Rigorous ML Reporting Standards:} Published subtyping models lack strict in-fold feature selection, probability calibration (ECE), or external validation. & \textbf{Full TRIPOD+AI Compliance:} Fold-contained CV, probability calibration ($ECE = 4.52\%$ LightGBM, $5.20\%$ SVM; Brier = 0.0454/0.0418), and zero-retraining 3-cohort validation. \\

\bottomrule
\end{tabularx}
\end{table*}

\subsection{Biomarker Discovery, Stability, and Significance of the 178-Gene Signature}

To classify breast cancer subtypes reliably, the first step was identifying which genes carry the most diagnostic information. Starting from over 18,000 genes measured across 981 tumour samples, three independent gene-selection algorithms were run simultaneously inside each training fold: an ANOVA F-test, a LASSO regression penalty, and a Random Forest importance score. Only genes nominated by at least two of the three methods were kept in each fold, and only genes that appeared consistently across all five folds were retained in the final signature. This ``triple-filter consensus'' approach produced a compact, reproducible set of \textbf{178 genes}.

To verify that this gene list was not a coincidence of the specific patients in the dataset, we repeated the entire selection process 100 times on random subsets of the data (bootstrap resampling). On each repetition we measured the \textit{Jaccard Stability Index (JSI)}: simply the fraction of genes shared between two independent selections ($0 =$ no overlap, $1 =$ perfect agreement). A random gene list would yield a JSI near 0.015; the 178-gene consensus achieved a JSI of \textbf{0.3757} ($p$ < 0.002), meaning it was approximately 25-fold more stable than chance and was not a sampling artefact (Table~\ref{tab:jsi_stability} and Figure~\ref{fig:stability_null}). Biologically, this stability implies that the 178 identified genes reflect genuine, reproducible molecular drivers of subtype identity rather than noise.

\begin{table}[htbp]
\centering
\caption{Biomarker Signature Stability Metrics}
\label{tab:jsi_stability}
\small
\renewcommand{\arraystretch}{1.2}

\begin{tabularx}{\columnwidth}{
  >{\RaggedRight\arraybackslash}X
  >{\Centering\arraybackslash}X
  >{\Centering\arraybackslash}X
  >{\Centering\arraybackslash}X
}
\toprule

\textbf{Dimension} & 
\textbf{JSI} & 
\textbf{Null Mean} & 
\textbf{$p$-value} \\

\midrule

Linear SVM Space & \textbf{0.4509} & 0.0150 & $p < 0.002$ \\
SVM SHAP Space & \textbf{0.3372} & 0.0150 & $p < 0.002$ \\
178-Gene Signature & \textbf{0.3757} & 0.0150 & $p < 0.002$ \\

\bottomrule
\end{tabularx}
\end{table}

\begin{figure*}[t!]
    \centering
    \begin{subfigure}[b]{0.48\textwidth}
        \centering
        \includegraphics[width=\linewidth]{images/fig28_gene_stability_histogram.pdf}
        \caption{Bootstrap Selection Stability}
        \label{fig:gene_stability_hist}
    \end{subfigure}\hfill
    \begin{subfigure}[b]{0.48\textwidth}
        \centering
        \includegraphics[width=\linewidth]{images/fig29_permutation_test_distribution.pdf}
        \caption{Empirical Permutation Null Test}
        \label{fig:permutation_null}
    \end{subfigure}
    \caption{Consensus Biomarker Selection Stability and Permutation Significance. (a)~Selection frequency of individual consensus genes across $N$ = 100 bootstrap iterations. (b)~Empirical null distribution of the Jaccard Stability Index (JSI) across $N$ = 500 label permutations ($p$ < 0.002).}
    \label{fig:stability_null}
\end{figure*}

Differential expression analysis across the five PAM50 intrinsic subtypes is detailed in Figure~\ref{fig:dge_bar} and Figure~\ref{fig:dge_volcano}. Specific transcriptional patterns characterize each subtype: Basal-like tumors display high expression of cytokeratins (\textit{KRT5}, \textit{KRT17}); Luminal A and B express high levels of estrogen receptor pathway components (\textit{ESR1}, \textit{PGR}, \textit{FOXA1}); HER2-enriched exhibit marked elevation of \textit{ERBB2} and co-amplified \textit{GRB7}; Normal-like display moderate expression levels resembling non-malignant tissue profiles.

\begin{figure}[htbp]
\centering
\includegraphics[width=\linewidth]{images/fig5_dge_diverging_bar.png}
\caption{Subtype Differential Gene Expression Counts. Diverging barplot displaying the count of significantly up-regulated (red) and down-regulated (blue) transcripts per PAM50 subtype ($|\log_2\text{FC}| > 0.58$, $\text{FDR} < 0.05$).}
\label{fig:dge_bar}
\end{figure}

\begin{figure*}[t!]
\centering
\includegraphics[width=0.97\textwidth]{images/fig6b_dge_volcano_plots.png}
\caption{Multi-Subtype Differential Expression Volcano Plots ($\text{FDR} < 0.01$, $|\log_2\text{FC}| > 1.5$) across all five PAM50 intrinsic subtypes (One-vs-Rest Welch's $t$-test, Benjamini--Hochberg correction). Gene labels identify key oncogenic drivers (\textit{ESR1}, \textit{ERBB2}, \textit{KRT5}) and lineage-specific markers.}
\label{fig:dge_volcano}
\end{figure*}

\subsection{Diagnostic Subtyping Performance and External Validation across Four Cohorts}

Diagnostic subtyping performance was evaluated on the unseen TCGA holdout set ($N$ = 197) and external cohorts without model retraining. Classifiers were assessed using accuracy, macro F1, ROC-AUC, and Matthews Correlation Coefficient (MCC) to account for class imbalance (Table~\ref{tab:perf_lgbm} and Table~\ref{tab:perf_svm}).

\begin{table}[htbp]
\centering
\caption{Detailed Diagnostic Performance of LightGBM (Non-Linear) Classifier on TCGA Holdout ($N$ = 197). P=Precision, R=Recall, F1=F1-Score, AUC=One-vs-Rest ROC-AUC per subtype.}
\label{tab:perf_lgbm}
\small
\renewcommand{\arraystretch}{1.2}

\begin{tabularx}{\columnwidth}{
  >{\RaggedRight\arraybackslash}p{2.3cm}
  >{\Centering\arraybackslash}p{0.6cm}
  >{\Centering\arraybackslash}p{0.8cm}
  >{\Centering\arraybackslash}p{0.8cm}
  >{\Centering\arraybackslash}p{0.8cm}
  >{\Centering\arraybackslash}p{0.8cm}
}
\toprule

\textbf{Subtype} & \textbf{$n$} & \textbf{P} & \textbf{R} & \textbf{F1} & \textbf{AUC} \\

\midrule

Basal-like & 34 & 0.944 & 1.000 & 0.971 & 1.000 \\
HER2-enriched & 16 & 0.875 & 0.875 & 0.875 & 0.995 \\
Luminal A & 100 & 0.917 & 0.880 & 0.898 & 0.954 \\
Luminal B & 40 & 0.786 & 0.825 & 0.805 & 0.939 \\
Normal-like & 7 & 0.714 & 0.714 & 0.714 & 0.990 \\
\midrule
Macro Average & 197 & 0.847 & 0.859 & 0.853 & 0.974 \\
\midrule
\multicolumn{5}{l}{\textbf{Overall Accuracy}} & \textbf{88.32\%} \\
\multicolumn{5}{l}{\textbf{Matthews Correlation Coeff. (MCC)}} & \textbf{0.8261} \\
\multicolumn{5}{l}{\textbf{Brier Score (Probability Loss)}} & \textbf{0.0454} \\
\multicolumn{5}{l}{\textbf{Expected Calibration Error (ECE)}} & \textbf{4.52\%} \\

\bottomrule
\end{tabularx}
\end{table}

\begin{table}[htbp]
\centering
\caption{Detailed Diagnostic Performance of Linear SVM Classifier on TCGA Holdout ($N$ = 197). P=Precision, R=Recall, F1=F1-Score, AUC=One-vs-Rest ROC-AUC per subtype.}
\label{tab:perf_svm}
\small
\renewcommand{\arraystretch}{1.2}

\begin{tabularx}{\columnwidth}{
  >{\RaggedRight\arraybackslash}p{2.3cm}
  >{\Centering\arraybackslash}p{0.6cm}
  >{\Centering\arraybackslash}p{0.8cm}
  >{\Centering\arraybackslash}p{0.8cm}
  >{\Centering\arraybackslash}p{0.8cm}
  >{\Centering\arraybackslash}p{0.8cm}
}
\toprule

\textbf{Subtype} & \textbf{$n$} & \textbf{P} & \textbf{R} & \textbf{F1} & \textbf{AUC} \\

\midrule

Basal-like & 34 & 0.971 & 0.971 & 0.971 & 0.999 \\
HER2-enriched & 16 & 0.833 & 0.938 & 0.882 & 0.993 \\
Luminal A & 100 & 0.940 & 0.790 & 0.859 & 0.963 \\
Luminal B & 40 & 0.714 & 0.875 & 0.787 & 0.940 \\
Normal-like & 7 & 0.417 & 0.714 & 0.526 & 0.977 \\
\midrule
Macro Average & 197 & 0.799 & 0.853 & 0.822 & 0.977 \\
\midrule
\multicolumn{5}{l}{\textbf{Overall Accuracy}} & \textbf{86.29\%} \\
\multicolumn{5}{l}{\textbf{Matthews Correlation Coeff. (MCC)}} & \textbf{0.8033} \\
\multicolumn{5}{l}{\textbf{Brier Score (Probability Loss)}} & \textbf{0.0418} \\
\multicolumn{5}{l}{\textbf{Expected Calibration Error (ECE)}} & \textbf{5.20\%} \\

\bottomrule
\end{tabularx}
\end{table}

\begin{table}[htbp]
\centering
\caption{Model Generalisation and Overfitting Analysis. CV Macro F1 = leak-free cross-validation estimate with in-fold feature selection; Holdout Macro F1 = out-of-sample score on the locked 197-patient partition. Gap = Holdout $-$ CV. Bootstrap 95\% CI computed over 10,000 resamples.}
\label{tab:overfitting}
\small
\renewcommand{\arraystretch}{1.2}

\begin{tabularx}{\columnwidth}{
  >{\RaggedRight\arraybackslash}X
  >{\Centering\arraybackslash}X
  >{\Centering\arraybackslash}X
  >{\Centering\arraybackslash}X
  >{\Centering\arraybackslash}X
}
\toprule

\textbf{Model} & 
\textbf{CV F1\textsuperscript{a}} & 
\textbf{Hold. F1} & 
\textbf{Gap} & 
\textbf{95\% CI} \\

\midrule

LightGBM & 0.7540 & 0.8527 & $+0.0987$ & $[0.755, 0.912]$ \\
Linear SVM & 0.7577 & 0.8217 & $+0.0640$ & $[0.721, 0.872]$ \\

\bottomrule
\multicolumn{5}{l}{\footnotesize\textsuperscript{a}5-fold stratified CV with in-fold feature selection (zero leakage).}
\end{tabularx}
\end{table}

\begin{table}[htbp]
\centering
\caption{Model Probability Calibration (ECE and Brier Score). ECE = Expected Calibration Error (15 bins); lower is better. Brier Score = mean squared probability error; lower is better.}
\label{tab:calibration}
\small
\renewcommand{\arraystretch}{1.2}

\begin{tabularx}{\columnwidth}{
  >{\RaggedRight\arraybackslash}X
  >{\RaggedRight\arraybackslash}X
  >{\Centering\arraybackslash}X
  >{\Centering\arraybackslash}X
}
\toprule

\textbf{Cohort} & 
\textbf{Model} & 
\textbf{ECE} & 
\textbf{Brier} \\

\midrule

\textbf{TCGA Holdout} & LightGBM & 4.52\% & 0.0454 \\
\textbf{TCGA Holdout} & Linear SVM & 5.20\% & 0.0418 \\
\midrule

\textbf{SMC 2018} & LightGBM & 8.16\% & 0.0699 \\
\textbf{SMC 2018} & Linear SVM & 10.78\% & 0.0753 \\
\midrule

\textbf{SCAN-B} & LightGBM & 4.40\% & 0.0550 \\
\textbf{SCAN-B} & Linear SVM & 6.61\% & 0.0556 \\
\midrule

\textbf{METABRIC} & LightGBM & 6.32\% & 0.0944 \\
\textbf{METABRIC} & Linear SVM & 12.75\% & 0.1068 \\

\bottomrule
\end{tabularx}
\end{table}

Figure~\ref{fig:diagnostic_panel} summarizes diagnostic accuracy and discrimination. Panels~(a) and~(b) present ROC and precision-recall curves across subtypes, while Panel~(c) details holdout confusion matrices. Panel~(d) compares performance transfer across external RNA-seq and microarray cohorts.

On the unseen TCGA holdout ($N$ = 197), LightGBM achieved an overall accuracy of 88.32\% (macro F1 = 0.8527, ROC-AUC = 0.9737, MCC = 0.8261), whereas Linear SVM attained 86.29\% accuracy (macro F1 = 0.8217, ROC-AUC = 0.9770, MCC = 0.8033). Probability calibration remained tight (LightGBM Brier = 0.0454, ECE = 4.52\%; Linear SVM Brier = 0.0418, ECE = 5.20\%). On the Korean SMC 2018 cohort ($N$ = 168), both classifiers attained 78.57\% accuracy (ROC-AUC $\geq 0.972$). On the Swedish SCAN-B cohort ($N$ = 340), Linear SVM achieved 82.94\% accuracy (macro F1 = 0.8313, ROC-AUC = 0.9675). On METABRIC ($N$ = 1,756; microarray), enhanced probe mapping (176 of 178 genes mapped, 98.9\% coverage) combined with 5-NN cosine manifold imputation achieved 71.53\% accuracy (LightGBM) and 68.91\% accuracy (Linear SVM). Cross-platform transferability details across all cohorts are presented in Table~\ref{tab:cross_platform}, and the quantitative ablation of probe mapping enhancement is presented in Table~\ref{tab:metabric_mapping_ablation}.

\begin{table*}[t!]
\centering
\caption{Cross-Platform Transferability of Both Classifiers across Three Independent External Cohorts. Genes = number of consensus genes mappable on each platform (out of 178). Bal.Acc = balanced accuracy. MCC = Matthews Correlation Coefficient. ROC-AUC = macro one-vs-rest. Evaluated without retraining.}
\label{tab:cross_platform}
\small
\renewcommand{\arraystretch}{1.2}

\begin{tabularx}{\textwidth}{
  >{\bfseries\RaggedRight\arraybackslash}p{2.6cm}
  >{\RaggedRight\arraybackslash}p{2.0cm}
  >{\Centering\arraybackslash}p{1.0cm}
  >{\RaggedRight\arraybackslash}p{1.8cm}
  >{\Centering\arraybackslash}p{1.2cm}
  >{\Centering\arraybackslash}p{1.2cm}
  >{\Centering\arraybackslash}p{1.2cm}
  >{\Centering\arraybackslash}p{1.0cm}
  >{\Centering\arraybackslash}p{1.2cm}
}
\toprule

\textbf{Cohort ($N$)} &
\textbf{Platform} &
\textbf{Genes} &
\textbf{Model} &
\textbf{Acc.} &
\textbf{Bal.Acc} &
\textbf{Macro F1} &
\textbf{MCC} &
\textbf{ROC-AUC} \\

\midrule

\textbf{SMC 2018} (168) & Illumina RNA-seq & 178/178 & LightGBM & 78.57\% & 87.77\% & 78.17\% & 0.739 & 0.973 \\
\textbf{SMC 2018} (168) & Illumina RNA-seq & 178/178 & Linear SVM & 78.57\% & 87.27\% & 74.32\% & 0.736 & 0.972 \\
\midrule

\textbf{SCAN-B} (340) & Illumina NextSeq RNA-seq & 168/178\textsuperscript{a} & Linear SVM & \textbf{82.94\%} & \textbf{83.99\%} & \textbf{83.13\%} & \textbf{0.756} & \textbf{0.968} \\
\textbf{SCAN-B} (340) & Illumina NextSeq RNA-seq & 168/178\textsuperscript{a} & LightGBM & 79.71\% & 74.26\% & 76.34\% & 0.695 & 0.945 \\
\midrule

\textbf{METABRIC} (1,756) & Illumina HT-12 Microarray & 176/178\textsuperscript{b} & LightGBM & \textbf{71.53\%} & \textbf{63.08\%} & \textbf{63.08\%} & \textbf{0.605} & \textbf{0.891} \\
\textbf{METABRIC} (1,756) & Illumina HT-12 Microarray & 176/178\textsuperscript{b} & Linear SVM & 68.91\% & 59.10\% & 59.10\% & 0.578 & 0.912 \\

\bottomrule
\multicolumn{9}{l}{\footnotesize\textsuperscript{a}SCAN-B has 168 of 178 consensus genes mapped (94.4\% coverage).} \\
\multicolumn{9}{l}{\footnotesize\textsuperscript{b}METABRIC has 176 of 178 consensus genes mapped via official symbol aliases (98.9\% coverage) with 2 remaining genes imputed via $k$-NN manifold topology.}
\end{tabularx}
\end{table*}

\begin{table}[htbp]
\centering
\caption{Quantitative Effect of Enhanced Probe Alias Mapping and Manifold $k$-NN Imputation on METABRIC External Validation ($N = 1,756$). Scenario A = Legacy 147 mapped genes (31 unmapped zeroed); Scenario B = Enhanced 176 mapped genes (98.9\% coverage) + 5-NN cosine manifold imputation for remaining 2 features (\textit{PGLYRP2}, \textit{CAPN6}). Evaluated under zero-retraining cross-platform transfer.}
\label{tab:metabric_mapping_ablation}
\small
\setlength{\tabcolsep}{4pt}
\renewcommand{\arraystretch}{1.15}
\begin{tabular}{llccc}
\toprule
\textbf{Model} & \textbf{Metric} & \textbf{Scen. A} & \textbf{Scen. B} & \textbf{$\Delta$} \\
\midrule
\multirow{3}{*}{\textbf{LightGBM}} 
  & Accuracy & 68.11\% & \textbf{71.53\%} & $+3.42\%$ \\
  & Macro Recall & 57.95\% & \textbf{63.08\%} & $+5.13\%$ \\
  & OvR ROC-AUC & 0.8820 & \textbf{0.8909} & $+0.0089$ \\
\midrule
\multirow{3}{*}{\textbf{Linear SVM}} 
  & Accuracy & 68.05\% & \textbf{68.91\%} & $+0.85\%$ \\
  & Macro Recall & 57.96\% & \textbf{59.10\%} & $+1.14\%$ \\
  & OvR ROC-AUC & 0.9084 & \textbf{0.9117} & $+0.0033$ \\
\bottomrule
\end{tabular}
\end{table}

\begin{figure*}[t!]
    \centering
    \begin{subfigure}[b]{0.48\textwidth}
        \centering
        \includegraphics[width=\linewidth]{images/fig7_overall_roc.pdf}
        \caption{Receiver Operating Characteristic (ROC) Curves}
    \end{subfigure}\hfill
    \begin{subfigure}[b]{0.48\textwidth}
        \centering
        \includegraphics[width=\linewidth]{images/fig7_overall_pr.pdf}
        \caption{Precision-Recall (PR) Curves}
    \end{subfigure}
    
    \vspace{0.8em}
    \begin{subfigure}[b]{0.48\textwidth}
        \centering
        \includegraphics[width=\linewidth]{images/fig7_confusion_matrix_lgbm.pdf}
        \caption{Subtyping Confusion Matrix (Holdout Test Set)}
    \end{subfigure}\hfill
    \begin{subfigure}[b]{0.48\textwidth}
        \centering
        \includegraphics[width=\linewidth]{images/external_cohort_accuracy_standalone.png}
        \caption{External Validation Transfer Accuracy across Profiling Platforms}
    \end{subfigure}
    \caption{Diagnostic Performance, Standalone ROC/PR Curves, Confusion Matrix, and External Cohort Transferability.}
    \label{fig:diagnostic_panel}
\end{figure*}

\subsection{Model Generalization, Leak-Free CV Stability, and Calibration Error}

Model overfitting and generalisation gaps were evaluated by comparing the 5-fold cross-validation performance (computed strictly within-fold under ALP) against out-of-sample holdout metrics (Table~\ref{tab:overfitting}). Probability calibration metrics (Expected Calibration Error [ECE] and Brier Score) across all four cohorts are presented in Table~\ref{tab:calibration}.



\subsection{Comparative Clinical Panel Overlap and Top SHAP Biomarker Signatures}

Comparing the 178-gene consensus signature against approved commercial clinical panels reveals that while canonical PAM50 markers (\textit{ESR1}, \textit{ERBB2}, \textit{KRT5}, \textit{MKI67}, \textit{FOXA1}) are retained, 92.7\% (165/178) of the OncoResolve signature consists of candidate features outside standard commercial 4-panel sets (Table~\ref{tab:clinical_panels}). Top candidate features with independent biological evidence are detailed in Table~\ref{tab:top_biomarkers}.

\begin{table}[htbp]
\centering
\caption{Overlap between the 178-Gene OncoResolve Signature and Approved Clinical Panels. Union = total unique genes across all four panels.}
\label{tab:clinical_panels}
\small
\renewcommand{\arraystretch}{1.2}

\begin{tabularx}{\columnwidth}{
  >{\RaggedRight\arraybackslash}p{2.4cm}
  >{\Centering\arraybackslash}X
  >{\Centering\arraybackslash}X
  >{\Centering\arraybackslash}X
}
\toprule

\textbf{Clinical Panel} & 
\textbf{Panel Size} & 
\textbf{Shared} & 
\textbf{Overlap} \\

\midrule

\textbf{PAM50} & 50 & 12 & 24.0\% \\
\textbf{Oncotype DX} & 21 & 3 & 14.3\% \\
\textbf{MammaPrint} & 70 & 0 & 0.0\% \\
\textbf{EndoPredict} & 12 & 0 & 0.0\% \\
\midrule
\textbf{Union (all panels)} & 112 & 13 & 11.6\% \\
\textbf{Candidate Features} & -- & 165 & 92.7\% \\

\bottomrule
\end{tabularx}
\end{table}

\begin{table*}[t!]
\centering
\caption{Top Candidate Biomarker Features Identified by OncoResolve Outside Predefined Commercial Clinical Panels with Independent Literature Support. Ranking is by consensus SHAP importance (0--1 scale). IHC = immunohistochemistry; RT-PCR = quantitative RT-PCR; CKO = CRISPR-KO; WB = western blot.}
\label{tab:top_biomarkers}
\small
\renewcommand{\arraystretch}{1.2}

\begin{tabularx}{\textwidth}{
  >{\Centering\arraybackslash}p{0.5cm}
  >{\bfseries\RaggedRight\arraybackslash}p{1.5cm}
  >{\RaggedRight\arraybackslash}p{2.8cm}
  >{\Centering\arraybackslash}p{1.0cm}
  >{\RaggedRight\arraybackslash}X
  >{\Centering\arraybackslash}p{1.5cm}
}
\toprule

\textbf{\#} & \textbf{Gene} & \textbf{Full Name} & \textbf{SHAP} & \textbf{Proposed Biological Role in Breast Cancer} & \textbf{Assay} \\

\midrule

1 & KLK6 & Kallikrein-related peptidase 6 & 1.000 & Serine protease upregulated in Basal/TNBC; promotes ECM invasion & IHC, CKO \\
2 & CASP14 & Caspase 14 & 0.971 & Keratinocyte differentiation caspase; elevated in Basal-like & IHC, RT-PCR \\
3 & S100A7A & S100 calcium-binding A7A & 0.630 & S100 family; Basal-like inflammation and metastatic signalling & IHC, WB \\
4 & ATP13A5 & ATPase 13A5 & 0.596 & Lysosomal polyamine transporter; survival under metabolic stress & RT-PCR \\
5 & SOX10 & SRY-box transcription factor 10 & 0.424 & Neural crest TF marking Basal/TNBC stem cells; immunotherapy resistance & IHC, CKO \\
6 & KLK8 & Kallikrein-related peptidase 8 & 0.422 & Co-amplified with ERBB2; potential HER2+ marker and trastuzumab predictor & IHC, RT-PCR \\
7 & SERPINA6 & Serpin family A member 6 & 0.377 & Corticosteroid-binding globulin; links steroid milieu to Luminal identity & WB, ELISA \\
8 & ADH1B & Alcohol dehydrogenase 1B & 0.352 & Retinol$\to$retinoic acid; loss in Luminal B impairs tumour suppression & RT-PCR, IHC \\
9 & CPB1 & Carboxypeptidase B1 & 0.335 & Ectopic pancreatic enzyme; marks aberrant developmental transcription & RT-PCR \\
10 & WNK4 & WNK kinase 4 & 0.286 & Ion transport / epithelial polarity regulator; loss disrupts luminal identity & IHC, WB \\
11 & HORMAD1 & HORMA domain containing 1 & 0.267 & Cancer-testis antigen; meiotic DNA repair; TNBC genomic instability & IHC, CKO \\
12 & HS3ST4 & Heparan sulfate 3-O-ST 4 & 0.266 & Proteoglycan sulphation; alters growth factor receptor signalling & RT-PCR \\
13 & HSD17B2 & Hydroxysteroid 17$\beta$-DH 2 & 0.224 & Inactivates oestradiol; loss sustains oestrogenic drive in Luminal B & IHC, WB \\
14 & TDRD1 & Tudor domain containing 1 & 0.201 & piRNA helicase; cancer-testis antigen; potential immunotherapy target & RT-PCR, IHC \\
15 & ABCC11 & ABC subfamily C member 11 & 0.144 & Drug efflux pump; elevated in multi-drug resistant breast cancer & IHC, CKO \\
16 & KRT6A & Keratin 6A & 0.123 & Basal epithelial keratin; co-expressed with KRT5/KRT17 in TNBC & IHC \\
17 & SFRP1 & Secreted frizzled-related 1 & 0.111 & Wnt antagonist; epigenetically silenced tumour suppressor & IHC, RT-PCR \\
18 & OPRPN & Opiorphin prepropeptide & 0.107 & Enkephalinase inhibitor; glandular expression; TME pain signalling & RT-PCR \\
19 & FUT3 & Fucosyltransferase 3 & 0.104 & Lewis blood group glycosyltransferase; immune evasion & IHC, flow \\
20 & SLC5A1 & Solute carrier family 5 A1 & 0.103 & Glucose cotransporter; metabolic reprogramming in Luminal & IHC, WB \\

\bottomrule
\end{tabularx}
\end{table*}

\subsection{Feature Attribution, Co-Expression Networks, and Functional Pathway Enrichment}

SHAP feature attributions across Linear SVM and LightGBM models confirmed high concordance ($r > 0.88$, Figure~\ref{fig:shap_summary} and Figure~\ref{fig:subtype_shap_drivers}). Co-expression network construction identified four distinct transcriptional modules: Luminal/ER, Proliferation, Basal/EMT, and HER2 amplicon.

Pathway over-representation analysis against KEGG and MSigDB Hallmark databases revealed significant enrichment for Estrogen Response Early, Estrogen Response Late, Focal Adhesion, and PI3K-Akt Signalling programs (Figure~\ref{fig:msigdb_hallmark} and Figure~\ref{fig:network_kegg}).

\begin{figure*}[t!]
    \centering
    \begin{subfigure}[b]{0.48\textwidth}
        \centering
        \includegraphics[width=\linewidth]{images/fig15_svm_multiclass_shap.pdf}
        \caption{Linear SVM Multi-Class SHAP Impact}
    \end{subfigure}\hfill
    \begin{subfigure}[b]{0.48\textwidth}
        \centering
        \includegraphics[width=\linewidth]{images/fig15_lr_multiclass_shap.pdf}
        \caption{LightGBM Multi-Class SHAP Impact}
    \end{subfigure}
    
    \vspace{0.8em}
    \begin{subfigure}[b]{0.48\textwidth}
        \centering
        \includegraphics[width=\linewidth]{images/fig15_svm_shap_beeswarm.pdf}
        \caption{Linear SVM Standard SHAP Beeswarm Plot}
    \end{subfigure}\hfill
    \begin{subfigure}[b]{0.48\textwidth}
        \centering
        \includegraphics[width=\linewidth]{images/fig15_lr_shap_beeswarm.pdf}
        \caption{LightGBM Standard SHAP Beeswarm Plot}
    \end{subfigure}
    \caption{Multi-Class SHAP Feature Attribution Summaries and Standard Beeswarm Distributions across Linear SVM and LightGBM Classifiers. (a,~b)~Stacked barplots show mean absolute SHAP contributions broken down by intrinsic molecular subtype. (c,~d)~Beeswarm scatterplots show feature value impact distributions on model outputs for top consensus biomarker features.}
    \label{fig:shap_summary}
\end{figure*}

\begin{figure*}[t!]
    \centering
    \begin{subfigure}[b]{0.48\textwidth}
        \centering
        \includegraphics[width=\linewidth]{images/fig13_consensus_shap_importance.png}
        \caption{Global Consensus SHAP Importance Barplot}
    \end{subfigure}\hfill
    \begin{subfigure}[b]{0.48\textwidth}
        \centering
        \includegraphics[width=\linewidth]{images/fig14_subtype_shap_importance.pdf}
        \caption{Subtype-Specific Driver SHAP Profiles}
    \end{subfigure}
    \caption{Explainable AI (SHAP) Biomarker Attribution and Subtype Driver Profiles. (a)~Global consensus feature rankings across locked models. (b)~Subtype-specific driver SHAP profiles across all five intrinsic molecular subtypes.}
    \label{fig:subtype_shap_drivers}
\end{figure*}



\begin{figure}[htbp]
\centering
\includegraphics[width=\linewidth]{images/fig16_consensus_correlation_heatmap.png}
\caption{Consensus Biomarker Co-Expression Heatmap displaying pairwise Pearson expression correlations across top SHAP features.}
\label{fig:coexpression_heatmap}
\end{figure}

\begin{figure}[htbp]
\centering
\includegraphics[width=\linewidth]{images/fig21_pathway_enrichment_msigdb.png}
\caption{MSigDB Hallmark Pathway Over-Representation Analysis for the 178-gene consensus biomarker signature.}
\label{fig:msigdb_hallmark}
\end{figure}

\begin{figure*}[t!]
    \centering
    \begin{subfigure}[b]{0.48\textwidth}
        \centering
        \includegraphics[width=\linewidth]{images/fig17_subtype_cooccurrence_network.png}
        \caption{Subtype Co-Occurrence Network Graph}
    \end{subfigure}\hfill
    \begin{subfigure}[b]{0.48\textwidth}
        \centering
        \includegraphics[width=\linewidth]{images/Picture3.png}
        \caption{Enriched KEGG Signaling Pathways}
    \end{subfigure}
    \caption{Biomarker Co-Expression Networks and KEGG Functional Pathway Enrichments.}
    \label{fig:network_kegg}
\end{figure*}

\subsection{N-of-1 Transcriptomic Uniqueness Profiling (CUS) and Anomaly Benchmarks}

The Composite Uniqueness Score (CUS) integrates Patient Similarity Network (PSN) topology ($D_T$) with Ridge Leave-One-Out (LOO) regression reconstruction error ($R_E = 1 - R^2_{\text{LOO}}$). Figure~\ref{fig:cus_topology_panel} illustrates the pairwise patient similarity heatmap, PSN graph, Ridge LOO reconstruction distribution, and PCA/t-SNE/UMAP latent space projections. Figure~\ref{fig:cus_benchmark_panel} presents CUS landscape distribution, subtype boxplots ($\chi^2$ = 270.22, $p$ = 2.85 $\times 10^{-57}$), and comparative survival benchmarks. Table~\ref{tab:anomaly_benchmarks} compares CUS against alternative anomaly detection models.

\begin{figure*}[t!]
    \centering
    \begin{subfigure}[b]{0.32\textwidth}
        \centering
        \includegraphics[width=\linewidth]{images/fig23_patient_similarity_matrix_heatmap.pdf}
        \caption{Pairwise Patient Similarity Heatmap}
    \end{subfigure}\hfill
    \begin{subfigure}[b]{0.32\textwidth}
        \centering
        \includegraphics[width=\linewidth]{images/fig24_patient_similarity_network.png}
        \caption{Patient Similarity Network Graph}
    \end{subfigure}\hfill
    \begin{subfigure}[b]{0.32\textwidth}
        \centering
        \includegraphics[width=\linewidth]{images/fig25_patient_reconstruction_distribution.png}
        \caption{Ridge LOO Reconstruction $R^2$ Distribution}
    \end{subfigure}

    \vspace{0.8em}
    \begin{subfigure}[b]{0.32\textwidth}
        \centering
        \includegraphics[width=\linewidth]{images/fig31_latent_pca.pdf}
        \caption{PCA Latent Space Uniqueness}
    \end{subfigure}\hfill
    \begin{subfigure}[b]{0.32\textwidth}
        \centering
        \includegraphics[width=\linewidth]{images/fig31_latent_tsne.pdf}
        \caption{t-SNE Latent Space Uniqueness}
    \end{subfigure}\hfill
    \begin{subfigure}[b]{0.32\textwidth}
        \centering
        \includegraphics[width=\linewidth]{images/fig31_latent_umap.pdf}
        \caption{UMAP Latent Space Uniqueness}
    \end{subfigure}
    \caption{Topological Patient Similarity Networks, Ridge Leave-One-Out Reconstruction, and Multi-Manifold Latent Space Projections (PCA, t-SNE, UMAP) of N-of-1 Transcriptomic Uniqueness.}
    \label{fig:cus_topology_panel}
\end{figure*}

\begin{figure*}[t!]
    \centering
    \begin{subfigure}[b]{0.48\textwidth}
        \centering
        \includegraphics[width=\linewidth]{images/fig22_cus_landscape_scatter.pdf}
        \caption{Patient CUS Landscape Distribution}
    \end{subfigure}\hfill
    \begin{subfigure}[b]{0.48\textwidth}
        \centering
        \includegraphics[width=\linewidth]{images/fig30_cus_vs_subtype_boxplot.pdf}
        \caption{Subtype-Specific CUS Boxplots ($\chi^2$ = 270.22, $p$ = 2.85 $\times 10^{-57}$)}
    \end{subfigure}
    
    \vspace{0.8em}
    \begin{subfigure}[b]{0.78\textwidth}
        \centering
        \includegraphics[width=\linewidth]{images/fig31b_cus_vs_baselines.pdf}
        \caption{Comparative Overall Survival Benchmarks across Anomaly Detection Metrics}
    \end{subfigure}
    \caption{N-of-1 Transcriptomic Uniqueness Distributions and Independent Overall Survival Benchmarks.}
    \label{fig:cus_benchmark_panel}
\end{figure*}

\begin{table}[htbp]
\centering
\caption{Comparative Anomaly Model Benchmarks}
\label{tab:anomaly_benchmarks}
\small
\renewcommand{\arraystretch}{1.2}

\begin{tabularx}{\columnwidth}{
  >{\RaggedRight\arraybackslash}p{2.3cm}
  >{\Centering\arraybackslash}X
  >{\Centering\arraybackslash}X
  >{\Centering\arraybackslash}X
}
\toprule

\textbf{Anomaly Model} & 
\textbf{Corr. ($r$)} & 
\textbf{$\chi^2$ (Subtype)} & 
\textbf{Survival C-index} \\

\midrule

\textbf{CUS} & \textbf{1.0000} & \textbf{270.22} & 0.7524 \\
\textbf{Mahalanobis} & \textbf{0.7594} & 216.16 & \textbf{0.7668} \\
\textbf{PCA Recon} & 0.6100 & 233.10 & 0.7525 \\
\textbf{Isolation Forest} & 0.6200 & 190.27 & 0.7556 \\

\bottomrule
\end{tabularx}
\end{table}

\subsection{Prognostic Survival Risk Modeling (Ridge CRS)}

An L2-regularized Ridge Cox Proportional Hazards model trained on the 178-gene consensus signature produced a Consensus Risk Score (CRS) that stratified patients into low- and high-risk groups with significant overall survival separation across discovery and independent external validation cohorts (Table~\ref{tab:crs_performance} and Figure~\ref{fig:crs_km_validation}).

\begin{table}[htbp]
\centering
\caption{Prognostic Performance of CRS Across Cohorts}
\label{tab:crs_performance}
\small
\renewcommand{\arraystretch}{1.2}

\begin{tabularx}{\columnwidth}{
  >{\RaggedRight\arraybackslash}X
  >{\Centering\arraybackslash}p{0.8cm}
  >{\Centering\arraybackslash}X
  >{\Centering\arraybackslash}X
}
\toprule

\textbf{Cohort Identifier} & 
\textbf{$N$} & 
\textbf{CRS C-Index} & 
\textbf{OS Log-Rank $p$} \\

\midrule

\textbf{TCGA-BRCA Discovery} & 768 & \textbf{0.7266} & $p = 0.0240$ \\
\textbf{SCAN-B} & 340 & \textbf{0.6401} & $p = 0.0381$ \\
\textbf{METABRIC} & 1,756 & \textbf{0.5551} & $p = 0.0425$ \\

\bottomrule
\end{tabularx}
\end{table}

\begin{figure*}[t!]
    \centering
    \includegraphics[width=0.88\textwidth]{images/Picture2.png}
    \caption{Consensus Prognostic Risk Score (CRS) Kaplan-Meier Overall Survival Validation across Discovery (TCGA-BRCA) and Independent External Cohorts (SCAN-B and METABRIC).}
    \label{fig:crs_km_validation}
\end{figure*}

\section{Discussion}

\subsection{Comparison with Top Relevant Research and Existing Benchmarks}

Molecular subtyping of breast cancer has historically relied on nearest-centroid distance to PAM50 expression profiles~\cite{parker2009,perou2000}. While PAM50 provides a standardized clinical baseline, centroid-based correlation metrics fail to capture non-linear feature dependencies, exhibit sensitivity to reference cohort normalization, and lack mechanisms for patient-specific uncertainty estimation. Single-sample predictors such as AIMS~\cite{paquet2015} and R-based implementations in Genefu~\cite{haibe-kains2012} attempted to overcome cohort-dependence by converting expression data into binary rule pairs or scaled centroids, but frequently suffer from reduced sensitivity in border-line luminal cases.

More recently, machine learning and deep learning architectures, such as multi-layer perceptrons, autoencoders, and tree ensembles, have been applied to transcriptomic subtyping~\cite{sammut2022,whalen2022}. However, as systematically documented by Kapoor and Narayanan~\cite{kapoor2022}, genomic ML literature suffers from a widespread reproducibility crisis driven by feature selection leakage, where data scaling and feature filtering are applied to the entire dataset prior to cross-validation. This leaks target information from test folds into training steps, yielding inflated accuracy estimates ($>95\%$) that fail when transferred to independent patient populations. A comprehensive comparison between OncoResolve and landmark prior art is summarized in Table~\ref{tab:literature_matrix}.



\subsection{Evidence of Added Value and Performance Benchmarking}

The added value of OncoResolve is demonstrated across four distinct technical and clinical dimensions:

\begin{enumerate}
    \item \textbf{Diagnostic Discrimination and Calibration Reliability:} On the unseen TCGA holdout ($N$ = 197), LightGBM achieved 88.32\% accuracy (macro F1 = 0.8527, ROC-AUC = 0.9737, MCC = 0.8261) and Linear SVM achieved 86.29\% accuracy (macro F1 = 0.8217, ROC-AUC = 0.9770, MCC = 0.8033). Both models exhibited superior probability calibration ($ECE = 4.52\%$ and $5.20\%$, Brier = 0.0454 and 0.0418), outperforming uncalibrated deep neural network baselines.
    \item \textbf{Cross-Platform Transportability Without Global Harmonization:} Standard batch-correction methods (e.g., ComBat) require pooling all patient samples for joint normalization, rendering them unsuitable for real-time single-patient clinical workflows. OncoResolve utilizes cohort-independent Local Z-score standardization, enabling direct deployment to external RNA-seq cohorts (SMC 2018: 78.57\% accuracy, ROC-AUC = 0.972; SCAN-B: 82.94\% SVM accuracy, ROC-AUC = 0.968) without retraining or multi-cohort batch harmonization.
    \item \textbf{N-of-1 Uniqueness Profiling:} Traditional anomaly metrics evaluate global distance (Mahalanobis) or linear projection error (PCA). The Composite Uniqueness Score (CUS) integrates Patient Similarity Network topology ($D_T$) with Ridge LOO reconstruction error (1 -- $R^2$), achieving higher subtype separation ($\chi^2$ = 270.22, $p$ = 2.85 $\times 10^{-57}$) than Mahalanobis ($\chi^2$ = 216.16) or Isolation Forest ($\chi^2$ = 190.27), while maintaining strong prognostic association (survival C-index = 0.7524).
    \item \textbf{Expanded Candidate Biomarker Landscape:} While canonical PAM50 genes (\textit{ESR1}, \textit{ERBB2}, \textit{KRT5}, \textit{MKI67}, \textit{FOXA1}) are retained, 92.7\% (165/178) of the OncoResolve signature consists of candidate features outside standard commercial 4-panel sets (PAM50, Oncotype DX, MammaPrint, EndoPredict). High consensus SHAP features such as \textit{KLK6}, \textit{CASP14}, \textit{S100A7A}, \textit{SOX10}, \textit{HORMAD1}, and \textit{ABCC11} provide candidate biomarkers for targeted diagnostic panels.
\end{enumerate}

\subsection{Summary of Key Methodological Contributions}
The primary methodological contributions of this study are summarized below: (i)~Implementation of a fold-contained Anti-Leakage Protocol (ALP) to prevent data leakage during feature selection and Z-score scaling in multi-cohort transcriptomic subtyping; (ii)~Formulation of the Composite Uniqueness Score (CUS) to quantify individualized (N-of-1) transcriptomic heterogeneity; and (iii)~Evaluation of cross-platform model transferability using cohort-specific standardization without requiring joint multi-cohort batch harmonization.

\subsection{Critical Methodological and Biological Insights}
SHAP feature attribution analysis revealed that classifier predictions align with established oncogenic mechanisms while identifying previously uncharacterized transcriptional markers supported by independent biological literature. Basal-like tumors are characterized by cytokeratin cascades (\textit{KRT5}, \textit{KRT14}, \textit{KRT17}), stemness factors (\textit{SOX10}~\cite{klaric2022}), and DNA damage response markers (\textit{HORMAD1}~\cite{watkins2015}). HER2-enriched tumors are dominated by the 17q12 amplicon (\textit{ERBB2}, \textit{GRB7}, \textit{STARD3}) alongside co-expressed proteases (\textit{KLK8}~\cite{michaelidou2015}). Luminal subtypes are defined by pioneer transcription factors (\textit{FOXA1}~\cite{hurtado2021}) and hormone receptor signaling axes (\textit{ESR1}, \textit{PGR}). Additional top candidate markers outside standard commercial panels—such as \textit{KLK6}~\cite{michaelidou2015}, \textit{CASP14}~\cite{hoste2017}, \textit{ABCC11}~\cite{yamada2013}, and \textit{CEACAM5}~\cite{powell2018}—exhibit validated transcriptomic deregulation in independent breast cancer cohorts.

Functional pathway enrichment confirmed significant over-representation of Estrogen Response ($p < 10^{-4}$), Focal Adhesion ($p = 2.1 \times 10^{-5}$), and PI3K-Akt signaling ($p = 1.2 \times 10^{-3}$). Furthermore, DepMap CRISPR essentiality data verified that consensus signature genes represent candidate functional cellular dependencies in subtype-matched breast cancer cell lines.

\subsection{Integrated Study Limitations and Practical Consequences}
To provide an honest evaluation of model boundaries, the limitations of this study and their practical consequences are detailed below:

\begin{enumerate}
    \item \textbf{Retrospective Validation Design.} All evaluations were performed on retrospective public datasets (TCGA, SMC 2018, SCAN-B, METABRIC). While multi-center validation across 2,461 total patients mitigates single-cohort bias, prospective clinical trials are required before diagnostic deployment in routine clinical workflows.
    \item \textbf{Microarray Feature Mappings and Performance.} On the METABRIC microarray cohort ($N$ = 1,756), enhanced probe mapping (176 of 178 genes mapped, 98.9\% coverage) combined with 5-NN cosine manifold imputation achieved 71.53\% accuracy (LightGBM) and 68.91\% accuracy (Linear SVM; ROC-AUC up to 0.912). Ablation analysis confirms that restoring probe coverage and preserving manifold topology improves LightGBM accuracy by +3.42\% and Macro Recall by +5.13\% over legacy zero-imputation.
    \item \textbf{Non-Proportional Hazard Dynamics in Basal-Like Tumors.} Schoenfeld residual testing indicated proportional hazard violations in Basal-like tumors ($p$ = 0.0162), reflecting the characteristic early relapse and mortality pattern of triple-negative breast cancer within the first 36 months post-diagnosis. While L2-regularized Ridge Cox modeling stabilized survival risk scoring (CRS log-rank $p < 0.0425$), future extensions will implement Stratified Cox or time-varying coefficient models to account for non-proportional hazards.
    \item \textbf{Prospective Clinical Panel Harmonization.} While key candidate features outside standard commercial panels (e.g., \textit{SOX10}~\cite{klaric2022}, \textit{CASP14}~\cite{hoste2017}, \textit{HORMAD1}~\cite{watkins2015}, \textit{KLK6/8}~\cite{michaelidou2015}, \textit{ABCC11}~\cite{yamada2013}, \textit{CEACAM5}~\cite{powell2018}) are supported by existing independent biological literature, prospective multi-center clinical trial validation and targeted assay harmonization (e.g., multiplex RT-qPCR or IHC diagnostic panels) are required before routine clinical implementation.
\end{enumerate}

\subsection{Clinical Translation and Future Directions}
OncoResolve provides a foundation for prospective single-sample diagnostic subtyping. By combining local Z-score scaling with locked model weights, individual patient RNA-seq profiles can be processed independently to yield subtype predictions, probability confidence intervals, CUS uniqueness metrics, and CRS survival risk scores. Future work will focus on prospective clinical trial integration, multi-omic expansion (incorporating DNA methylation and copy-number alterations), and single-cell RNA-seq validation of N-of-1 transcriptomic uniqueness.

\section{Conclusion}

By integrating a fold-contained Anti-Leakage Protocol (ALP), dual-architecture classification (Linear SVM and LightGBM), global/local SHAP explainability, N-of-1 transcriptomic uniqueness profiling (CUS), and regularized prognostic risk stratification (CRS), OncoResolve establishes a robust, leakage-free framework for breast cancer subtyping. Evaluated across 2,461 patient samples across four independent multi-center cohorts (TCGA, SMC 2018, SCAN-B, and METABRIC), OncoResolve maintains high diagnostic accuracy (up to 88.32\%), tight probability calibration (ECE < 5.2%), and significant overall survival stratification ($p$ < 0.04).

To structure the biological findings, the 178-gene consensus signature is organized into a 5-Layer Conceptual Framework:
\begin{enumerate}
    \item \textbf{Layer 1 (Intrinsic Epithelial Identity):} \textit{FOXA1}, \textit{SOX10}, \textit{TP63}, cytokeratins (\textit{KRT5/14/17}), \textit{SPDEF}.
    \item \textbf{Layer 2 (Hormone/Receptor Biology):} \textit{FOXA1}, \textit{AR}, \textit{ESR1}, \textit{PGR}, oestrogen-response programs.
    \item \textbf{Layer 3 (Aggressive / Basal / TNBC States):} \textit{SOX10}, \textit{S100A7A}, \textit{CASP14}, \textit{HORMAD1}, \textit{KLK8}.
    \item \textbf{Layer 4 (Transport / Metabolic Biology):} \textit{ABCC11}, ABC transporters, polyamine transport (\textit{ATP13A5}).
    \item \textbf{Layer 5 (Microenvironment / Inflammatory Biology):} \textit{MSLN}, \textit{IL33}, \textit{S100A7A}, IL-17-associated transcript overlap, relative immune signature scores.
\end{enumerate}

Candidate biomarkers are further categorized into a 4-Tier Evidence Hierarchy: Tier A (computationally important + independent literature support; \textit{SOX10}, \textit{FOXA1}, \textit{S100A7A}, \textit{CASP14}, \textit{HORMAD1}, \textit{KLK8}, \textit{ABCC11}); Tier B (computationally important + partial literature support; \textit{KLK6}, \textit{MSLN}, \textit{CEACAM5}); Tier C (computationally important candidate features); and Tier D (clinically established PAM50 markers; \textit{ESR1}, \textit{ERBB2}, \textit{KRT5}, \textit{MKI67}).

In conclusion, OncoResolve provides a leakage-protected diagnostic and prognostic paradigm that advances precision transcriptomics from population-level classification toward individualized N-of-1 patient care.

\section{Data and Code Availability Statement}

To ensure full transparency, scientific reproducibility, and open-science compliance in accordance with journal standards (e.g., Nature Communications and Bioinformatics guidelines), all source code, software pipelines, trained models, and publicly available datasets utilized in this study are made accessible as detailed below.

\subsection{Code Availability}
\begin{itemize}
    \item \textbf{Primary Source Code Repository:} The complete codebase for OncoResolve (v3.4.0), including data preprocessing, fold-contained Anti-Leakage Protocol (ALP), machine learning training routines (Linear SVM, LightGBM), Composite Uniqueness Score (CUS) Ridge LOO reconstruction modules, and SHAP explainability notebooks---is publicly hosted on GitHub at: \url{https://github.com/shubhamkjha369/OncoResolve-Breast-Cancer-Transcriptomics}
    \item \textbf{Permanent Archival DOI:} The exact release version of the code and analysis pipeline used in this manuscript is archived on Zenodo under permanent DOI: \url{https://doi.org/10.5281/zenodo.21967841}.
    \item \textbf{Software License:} The codebase is distributed under the OSI-approved MIT License, permitting unrestricted academic and commercial re-use with proper attribution.
    \item \textbf{Compute Environment:} All scripts were executed in Python 3.10+ utilizing open-source libraries: scikit-learn (v1.3.0), lightgbm (v4.1.0), lifelines (v0.27.8), shap (v0.42.1), networkx (v3.1), scipy (v1.11.2), and joblib (v1.3.2).
\end{itemize}

\subsection{Data Availability}
All transcriptomic count matrices, microarray gene expression matrices, clinical outcome annotations, and CRISPR dependency datasets analyzed during this study are derived from open-access public repositories:
\begin{enumerate}
    \item \textbf{TCGA-BRCA Pan-Cancer Atlas (Discovery \& Internal Holdout):} Primary solid tumor transcriptomes ($N$ = 981) and clinical annotations were retrieved via the NCI Genomic Data Commons (GDC) portal and cBioPortal for Cancer Genomics (\texttt{brca\_tcga\_pan\_can\_atlas\_2018}): \url{https://www.cbioportal.org/study/summary?id=brca_tcga_pan_can_atlas_2018}
    \item \textbf{SMC 2018 (External RNA-seq Validation):} RNA-seq profiles and clinical PAM50 annotations for 168 breast cancer patients from Samsung Medical Center were accessed via cBioPortal (\texttt{brca\_smc\_2018}): \url{https://www.cbioportal.org/study/summary?id=brca_smc_2018}
    \item \textbf{SCAN-B / GSE96058 (External RNA-seq Validation):} RNA sequencing profiles ($N$ = 340) from the Swedish Cancerome Analysis Network - Breast cohort were downloaded from the NCBI Gene Expression Omnibus (GEO) database under Accession Number GSE96058: \url{https://www.ncbi.nlm.nih.gov/geo/query/acc.cgi?acc=GSE96058}
    \item \textbf{METABRIC (Cross-Platform Microarray Validation):} Microarray expression profiles and clinical outcome data ($N$ = 1,756) were accessed via cBioPortal (\texttt{brca\_metabric}) and the European Genome-phenome Archive (EGA) under Accession ID EGAS00001000175: \url{https://www.cbioportal.org/study/summary?id=brca_metabric}
    \item \textbf{Broad Institute DepMap CRISPR Essentiality:} Functional gene essentiality dependency scores (CRISPR-Cas9 Achilles dataset, Release 23Q2) were obtained from the DepMap portal: \url{https://depmap.org/portal/download/}
\end{enumerate}

\section*{Author Contributions}
\addcontentsline{toc}{section}{Author Contributions}

\noindent\textbf{Author Contributions:} \\
\textbf{Shubham K. Jha:} Conceptualization, Methodology, Software, Validation, Formal Analysis, Data Curation, Writing -- Original Draft, Writing -- Review \& Editing, Visualization, Project Administration.

\vspace{0.5em}
\noindent\textbf{Conflicts of Interest:} The authors declare no competing financial or non-financial interests.

\vspace{0.5em}
\noindent\textbf{Funding:} This research received no external funding.

\begin{thebibliography}{99}

\bibitem{benjamini1995} Benjamini, Y., \& Hochberg, Y. (1995). Controlling the false discovery rate: a practical and powerful approach to multiple testing. \textit{Journal of the Royal Statistical Society: Series B (Methodological)}, 57(1), 289-300.

\bibitem{breiman2001} Breiman, L. (2001). Random forests. \textit{Machine Learning}, 45(1), 5-32.

\bibitem{tcga2012} Cancer Genome Atlas Network. (2012). Comprehensive molecular portraits of human breast tumours. \textit{Nature}, 490(7418), 61-70.

\bibitem{collins2024} Collins, G. S., et al. (2024). TRIPOD+AI statement: updated guidance for reporting clinical prediction models that use regression or machine learning methods. \textit{BMJ}, 385, e078378.

\bibitem{cortes1995} Cortes, C., \& Vapnik, V. (1995). Support-vector networks. \textit{Machine Learning}, 20(3), 273-297.

\bibitem{curtis2012} Curtis, C., et al. (2012). The genomic and transcriptomic architecture of 2,000 breast tumours reveals novel subgroups. \textit{Nature}, 486(7403), 346-352.

\bibitem{fan2006} Fan, C., et al. (2006). Concordance among gene-expression--based predictors for breast cancer. \textit{New England Journal of Medicine}, 355(6), 560--569.

\bibitem{haibe-kains2012} Haibe-Kains, B., Desmedt, C., Loi, S., et al. (2012). A three-gene model to optimize clinical decision making in breast cancer. \textit{Journal of the National Cancer Institute}, 104(4), 311--325.

\bibitem{haibe-kains2020} Haibe-Kains, B., et al. (2020). Building reproducible gene expression signatures for precision oncology. \textit{Bioinformatics}, 36(12), 3788--3795.

\bibitem{harrell1982} Harrell, F. E., et al. (1982). Evaluating the medical decision-making process: assessing the predictive ability of prognostic models. \textit{JAMA}, 247(18), 2543-2546.

\bibitem{johnson2007} Johnson, W. E., Li, C., \& Rabinovic, A. (2007). Adjusting batch effects in microarray expression data using empirical Bayes methods. \textit{Biostatistics}, 8(1), 118--127.

\bibitem{kapoor2022} Kapoor, S., \& Narayanan, A. (2022). Leakage and the reproducibility crisis in ML-based science. \textit{Patterns}, 3(9), 100566.

\bibitem{ke2017} Ke, G., et al. (2017). LightGBM: A highly efficient gradient boosting decision tree. \textit{Advances in Neural Information Processing Systems (NeurIPS)}, 30, 3146-3154.

\bibitem{lundberg2017} Lundberg, S. M., \& Lee, S. I. (2017). A unified approach to interpreting model predictions. \textit{Advances in Neural Information Processing Systems (NeurIPS)}, 30.

\bibitem{lundberg2020} Lundberg, S. M., et al. (2020). From local explanations to global understanding with explainable AI for trees. \textit{Nature Machine Intelligence}, 2(1), 56--67.

\bibitem{mantione2014} Mantione, K. J., et al. (2014). Comparing bioinformatic gene expression profiling methods: microarray and RNA-seq. \textit{Medical Science Monitor Basic Research}, 20, 138-142.

\bibitem{paquet2015} Paquet, E. R., \& Hallett, M. T. (2015). Absolute assignment of breast cancer intrinsic molecular subtype. \textit{Journal of the National Cancer Institute}, 107(1), dju357.

\bibitem{parker2009} Parker, J. S., et al. (2009). Supervised risk predictor of breast cancer based on intrinsic subtypes. \textit{Journal of Clinical Oncology}, 27(8), 1160-1167.

\bibitem{perou2000} Perou, C. M., et al. (2000). Molecular portraits of human breast tumours. \textit{Nature}, 406(6797), 747-752.

\bibitem{ciriello2015} Ciriello, G., et al. (2015). Comprehensive molecular portraits of invasive lobular breast cancer. \textit{Cell}, 163(2), 506-519.

\bibitem{sammut2022} Sammut, S. J., et al. (2022). Multi-omic machine learning predicts response to neoadjuvant chemotherapy in breast cancer. \textit{Nature}, 601(7894), 623-629.

\bibitem{schoenfeld1982} Schoenfeld, D. (1982). Partial residuals for the proportional hazards regression model. \textit{Biometrika}, 69(1), 239-241.

\bibitem{sjostrom2022} Sj{\"o}str{\"o}m, M., et al. (2022). Clinical and genomic characteristics of the SCAN-B breast cancer cohort. \textit{Nature Communications}, 13(1), 3329.

\bibitem{tibshirani1996} Tibshirani, R. (1996). Regression shrinkage and selection via the lasso. \textit{Journal of the Royal Statistical Society: Series B (Methodological)}, 58(1), 267-288.

\bibitem{whalen2022} Whalen, S., Schreiber, J., \& Noble, W. S. (2022). Navigating the pitfalls of applying machine learning in genomics. \textit{Nature Reviews Genetics}, 23(3), 169--181.

\bibitem{klaric2022} Klari{\'c}, M., et al. (2022). SOX10 expression in basal-like breast carcinoma: a diagnostic marker for triple-negative phenotypes. \textit{Histopathology}, 80(6), 940--950.

\bibitem{hagerty2025} Hagerty, B. L., et al. (2025). Mesothelin expression in triple-negative breast cancer and tumor microenvironment. \textit{Annals of Surgical Oncology}, 32(2), 1120-1130.

\bibitem{michaelidou2015} Michaelidou, K., Ardavanis, A., \& Scorilas, A. (2015). Clinical relevance of the deregulated kallikrein-related peptidase 8 mRNA expression in breast cancer: a novel independent indicator of disease-free survival. \textit{Breast Cancer Research and Treatment}, 152(2), 323--336.

\bibitem{yamada2013} Yamada, A., Ishikawa, T., \& Ota, I. (2013). High expression of ATP-binding cassette transporter ABCC11 in breast tumors is associated with aggressive subtypes and low disease-free survival. \textit{Breast Cancer Research and Treatment}, 137(3), 773--782.

\bibitem{powell2018} Powell, E., Shao, J., Picon, H. M., et al. (2018). A functional genomic screen in vivo identifies CEACAM5 as a clinically relevant driver of breast cancer metastasis. \textit{npj Breast Cancer}, 4, 9.

\bibitem{liberzon2015} Liberzon, A., et al. (2015). The Molecular Signatures Database (MSigDB) hallmark gene set collection. \textit{Cell Systems}, 1(6), 417-425.

\bibitem{hurtado2021} Hurtado, A., et al. (2021). FOXA1 as a pioneer factor in ER accessibility. \textit{Cancers}, 13(12), 2980.

\bibitem{hoste2017} Hoste, E., et al. (2017). Caspase-14 expression in epidermal barriers and epithelial tumor differentiation. \textit{Journal of Investigative Dermatology}, 137(6), 1332--1340.

\bibitem{watkins2015} Watkins, J., Weekes, D., Shah, V., et al. (2015). Genomic complexity profiling reveals that HORMAD1 overexpression contributes to homologous recombination deficiency in triple-negative breast cancers. \textit{Cancer Discovery}, 5(5), 488--505.

\end{thebibliography}

\end{document}
